\chapter{Хиггс-разрешённые опоры спектрального пакета пятнадцать}
\label{ch:v6-spectral-transition-higgs-resolved-support-gate}

\section{Постановка различающего теста}

Предыдущие гейты установили, что класс пятнадцать является пакетом
поколения, раскладывается как $12+3$ и не имеет выведенной энергии
колокализации. Следующий вопрос локален во внутреннем пространстве:
какие минимальные опоры возникают после выбора ненулевого направления
Хиггса и какие из них действительно соединены рёбрами конечного
оператора Дирака?

Работа ведётся на наблюдаемом пакете
\begin{equation}
 H_{15}=Q_L\oplus L_L\oplus u_R\oplus d_R\oplus e_R,
 \qquad \dim_{\mathbb C}H_{15}=15.
\end{equation}
Правый нейтрино отсутствует согласно замороженной архитектуре $H_{15}$.

\section{Два хиггсовских направления}

При $H\ne0$ определим
\begin{equation}
 P_H=\frac{HH^\dagger}{H^\dagger H},
 \qquad
 P_{\widetilde H}=\frac{\widetilde H\widetilde H^\dagger}{H^\dagger H},
 \qquad
 \widetilde H=i\sigma_2\overline H.
\end{equation}
Они удовлетворяют
\begin{equation}
 P_H+P_{\widetilde H}=I_2,
 \qquad P_HP_{\widetilde H}=0,
\end{equation}
и преобразуются ковариантно под $SU(2)$. Поэтому на $H_{15}$ возникают
четыре ортогональных проектора:
\begin{align}
 P_u&=(I_3\otimes P_{\widetilde H})|_{Q_L}\oplus I_3|_{u_R},\\
 P_d&=(I_3\otimes P_H)|_{Q_L}\oplus I_3|_{d_R},\\
 P_e&=P_H|_{L_L}\oplus I_1|_{e_R},\\
 P_\nu&=P_{\widetilde H}|_{L_L}.
\end{align}
Их ранги равны соответственно
\begin{equation}
 6,\quad6,\quad2,\quad1,
\end{equation}
а сумма равна единице на $H_{15}$.

Машинный аудит дал остаток суммы $3.15\cdot10^{-16}$,
остаток идемпотентности $2.39\cdot10^{-16}$ и максимальный остаток
$SU(2)$-ковариантности $1.21\cdot10^{-15}$. Цветные коммутаторы равны
нулю.

\section{Какие опоры соединены оператором Дирака}

Три допустимых ребра
\begin{equation}
 Q_L\leftrightarrow u_R,
 \qquad Q_L\leftrightarrow d_R,
 \qquad L_L\leftrightarrow e_R
\end{equation}
точно сохраняют $P_u,P_d,P_e$. Для любого ненулевого значения трёх
юкавских амплитуд соответствующие блоки являются дираковскими парами.
Структурный тест с произвольными ненулевыми числами $0.7,0.4,0.2$ дал
ранг оператора $14$ и одну нулевую моду. Числа использованы только для
проверки ранга и не являются предсказанием масс.

Единственный нуль лежит в $P_\nu H_{15}$. Причина не в малости
коэффициента, а в отсутствии правого конца ребра на $H_{15}$.
Следовательно,
\begin{equation}
 15=6_u+6_d+2_e+1_\nu
\end{equation}
является точным разложением опор, но не спектром выведенных масс.

\section{Тёплицев и Real-ledger}

Компонентные классы равны их комплексным рангам:
\begin{equation}
 6+6+2+1=15.
\end{equation}
После KO6-удвоения Real-ранги равны $12,12,4,2$. Нормированные веса
\begin{equation}
 \frac6{105}+\frac6{105}+\frac2{105}+\frac1{105}=\frac17
\end{equation}
точно восстанавливают вес исходного дефекта. Это подтверждает
совместимость разложения с прежним классом, но не делает каждый блок
наблюдаемой элементарной частицей.

\begin{theorem}[Хиггс-разрешение пакета поколения]
На области $H\ne0$ наблюдаемый пакет $H_{15}$ канонически и
gauge-ковариантно раскладывается на опоры рангов $6,6,2,1$.
Три первых опоры поддерживают дираковские рёбра степени один. Последняя
опора является односторонней нейтринной линией без степени-один массового
ребра.
\end{theorem}

\begin{nogocriterion}[Запрет преждевременной элементарности]
Нельзя объявлять разложение $6+6+2+1$ выводом четырёх физических масс.
Текущий родитель не фиксирует юкавские амплитуды, абсолютный масштаб
Хиггса и нейтринный оператор высшей степени. Кроме того, проектор
$P_\nu$ не определён при $H=0$.
\end{nogocriterion}

\section{Следующий гейт}

Гейт
\texttt{version6\_spectral\_transition\_neutrino\_line\_parent\_gate}
должен проверить, является ли $P_\nu(H)$ глобальным родительским
объектом или нормировка скрывает сингулярность. Отдельно следует
проверить ненормированный квадратичный носитель
$\widetilde H\widetilde H^\dagger$.

\begin{protocol}
\begin{enumerate}
 \item разложение $15=6+6+2+1$: \textbf{получено};
 \item gauge-ковариантность проекторов: \textbf{получена при $H\ne0$};
 \item ранг заряженного оператора: \textbf{$14$};
 \item число нейтринных нулей степени один: \textbf{один};
 \item значения заряженных масс: \textbf{не выведены};
 \item глобальность нейтринной линии через $H=0$: \textbf{открыта};
 \item физическое замыкание: \textbf{не достигнуто}.
\end{enumerate}
\end{protocol}

Машинный сертификат сохранён в
\path{s2t_v6_spectral_transition_higgs_resolved_support_gate.py} и
\path{s2t_v6_spectral_transition_higgs_resolved_support_gate_results.json}.
